\documentclass[10pt,a4paper]{amsart}
\usepackage[margin=19mm]{geometry}
\usepackage{amsmath,amssymb,mathtools}
\usepackage[scr=rsfs,cal=euler]{mathalfa}
\usepackage{xfrac}
\usepackage[unicode,hidelinks,bookmarks=false]{hyperref}
\newcommand{\ie}{i.e.\ }
\newcommand{\cf}{cf.\ }
\newcommand{\resp}{respectively\ }
\newcommand{\Subsection}{\subsection}
\providecommand{\nobreakdash}{\nobreak\hbox{-}}
\setlength{\emergencystretch}{2em}
\multlinegap0pt
\usepackage{amssymb}
\usepackage[shortlabels]{enumitem}
\newtheorem{proposition}{Proposition}
\usepackage{cleveref}
\crefname{proposition}{Proposition}{Propositions}
\newcommand{\Pe}{\mathrm{Pe}}
\newcommand{\QF}{\mathcal{Q}}
\newcommand{\R}{\mathbb{R}}
\newcommand{\abs}[1]{\left\vert #1 \right\vert}
\newcommand{\avg}[1]{\int_\Omega #1 \,\dx}
\newcommand{\cal}[1]{\mathcal{#1}}
\renewcommand{\div}{\divergence}%{\nabla \cdot }
\DeclareMathOperator{\divergence}{div}
\newcommand{\dx}{\mathrm{d} \vec{x}}
\newcommand{\invlap}{\Delta^{-1}}
\newcommand{\lb}{\mathcal{L}}
\newcommand{\norm}[1]{\left\| #1 \right\|}
\newcommand{\ocpval}{\mathcal{J}}
\renewcommand{\vec}[1]{\mathbf{#1}}
\title{Optimizing bounds for energy-constrained optimal cooling problems in two dimensions}
\author{Pedro Bl\"oss Braga, Giovanni Fantuzzi}
\date{Source: arXiv:2608.14334v1 (CC BY 4.0)}
\begin{document}
\maketitle

\noindent\emph{Source and licence.} Optimizing bounds for energy-constrained optimal cooling problems in two dimensions. Version arXiv:2608.14334v1 (CC BY 4.0), distributed by arXiv under the Creative Commons Attribution 4.0 International licence, http://creativecommons.org/licenses/by/4.0/, as recorded in the arXiv licence metadata for this exact version and shown on the abstract page https://arxiv.org/abs/2608.14334v1. Full licence notice: \texttt{licenses/arxiv-2608.14334-CC-BY-4.0.txt}.\par\smallskip

\noindent\textit{Editorial scope.} Counted unit: the article's proposition th:duality (source0.tex lines 291--302). The article's complete original proof of it is delivered with it (source0.tex lines 329--400, one \texttt{proof} environment of 72 lines, which the article places away from the statement). No mathematics is written, repaired or re-derived: the statement and the proof are the article's own lines, in the article's own notation, and no retained line is substituted or re-flowed.  Retained uncounted as the source-local context the proof needs: lines 176--181; lines 184--193; lines 272--274; lines 286--288.  Nothing was omitted from any retained span, so the package carries no omission record.  No retained line contains a citation, so the wrapper carries no bibliography and no bibliography entry.  No retained line contains a figure environment, an image-inclusion command or drawing code, so no image is delivered and the package declares no asset dependency.  Editorial formatting only, in addition: an AMS \texttt{amsart} preamble that restates the article's own declarations and macros used by the retained lines, namely the environments \texttt{proposition} on separate counters, together with the standard packages those lines need.  The retained environments print in this document as Proposition 1 in order of appearance, whereas the article prints the same items with the numbering of its own full document.  The complete native source of the article as distributed in the arXiv e-print for this exact version is bundled unchanged as \texttt{sources/207659/source0.tex}.
\begin{equation}\label{e:ade}
    \begin{cases}
        - \Delta T + \vec{u}\cdot \nabla T =  f &\text{on }\Omega,\\
        T=0 &\text{on }\partial\Omega.
    \end{cases}
\end{equation}

\begin{equation}\label{e:ocp}\tag{OCP}
    \ocpval(\Pe):=
    \min_{\substack{
        \vec{u}\in L^2(\Omega;\R^2) \\ 
        \div\vec{u}=0 \text{ on }\Omega\\ 
        \vec{u} \cdot \hat{\vec{n}} = 0  \text{ on }\partial\Omega\\
        \norm{\vec{u}}_{L^2} \leq \,\Pe
    }}\;
    \avg{\abs{\nabla T}^2}
\end{equation}

\begin{equation}\label{e:ade-weak}
    \avg{ \nabla \xi \cdot \nabla T + \Pe\,\xi \,\nabla^\perp \psi \cdot \nabla T - f\xi}= 0 \quad \forall \xi \in H^1_0(\Omega).
\end{equation}

\begin{equation}\label{e:qf}
    \mathcal{Q}\left(a,\xi\right)[\psi,\theta] := \avg{ a|\nabla \psi|^2 + \abs{\nabla\theta}^2 - \Pe\,\xi\, \nabla^\perp \psi \cdot \nabla \theta}.
\end{equation}

\begin{proposition}\label{th:duality}
    For every $\Pe\geq 0$, 
    the optimal value $\ocpval(\Pe)$ of problem \cref{e:ocp} satisfies $\ocpval(\Pe) \geq \lb(\Pe) \geq 0$, where
    \begin{equation}\label{e:ldp}\tag{LDP}
        \lb(\Pe) := \max_{\substack{
            a\in \R \\ 
            \xi \in H^1_0(\Omega) \\ 
            \QF\left(a,\xi\right) \succeq  0 
        }}\;
        \avg{f\xi - \frac14 \abs{\nabla \xi}^2} - a.
    \end{equation}
\end{proposition}

\begin{proof}[Proof of \cref{th:duality}]
    The temperature $T$ in \cref{e:ocp} is the unique function in $H^1_0(\Omega)$ satisfying \cref{e:ade-weak}. Viewing the test function $\xi$ in that equation as a Lagrange multiplier for the advection-diffusion equation \cref{e:ade}, and introducing a Lagrange multiplier $a\geq 0$ for the constraint $\norm{\nabla\psi}_{L^2}\leq 1$, we can write
    \begin{align*}
        \ocpval(\Pe)
        &= 
        \inf_{\substack{
        \psi \in \cal{H} \\ 
        T \in H^1_0(\Omega)
        }}\;
        \sup_{\substack{a \in \R \\ \xi \in H^1_0(\Omega) \\ a\geq 0}}
        \avg{\abs{\nabla T}^2 + a |\nabla \psi|^2 + (\nabla\xi - \Pe\,\xi\,\nabla^\perp\psi )\cdot \nabla T + f\xi} - a
        \\
        &\geq 
        \sup_{\substack{a \in \R \\ \xi \in H^1_0(\Omega) \\ a\geq 0}}
        \inf_{\substack{
        \psi \in \cal{H} \\ 
        T \in H^1_0(\Omega)
        }}\;
        \avg{\abs{\nabla T}^2 + a |\nabla \psi|^2 + (\nabla\xi - \Pe\,\xi\,\nabla^\perp\psi ) + f\xi} - a.
    \end{align*}
    The inequality on the second line is due to switching minimization and maximization.

    Next, we perform the inner minimization over $T$ for fixed $\psi$, $a$ and $\xi$. The optimal temperature $T^*$ solves the Euler--Lagrange equation
    $-2\Delta T^* - \Delta \xi + \Pe \, \nabla \xi \cdot \nabla^\perp \psi = 0$
    with boundary condition $T^* = 0$. Writing $\invlap$ for the inverse Laplacian with homogeneous Dirichlet boundary conditions, we obtain
    \begin{equation}\label{e:Topt}
        T^* = - \frac12\xi + \frac{\Pe}{2}\ \invlap(\nabla \xi \cdot \nabla^\perp \psi).
    \end{equation}
    To compute the value of the minimum, we multiply the Euler--Lagrange equation by $T^*$ and integrate by parts over $\Omega$ to find that
    \begin{equation}\label{e:el-1}
    \avg{ (\nabla\xi - \Pe\,\xi\,\nabla^\perp\psi )\cdot \nabla T^*} = -2 \avg{\abs{\nabla T^*}^2}.
    \end{equation}
    We then observe that
    \begin{equation}\label{e:el-2}
        \avg{\nabla \xi \cdot \nabla \invlap(\nabla \xi \cdot \nabla^\perp \psi)}
        = -\avg{ \xi \nabla \xi \cdot \nabla^\perp \psi}
        = 0,
    \end{equation}
    where the both equalities follows from integration by parts using the divergence-free nature of $\nabla^\perp \psi$ and the homogeneous Dirichlet boundary conditions of $\xi$. Using \cref{e:Topt}, \cref{e:el-1} and \cref{e:el-2}, we obtain
    \begin{align*}
        \ocpval(\Pe)
        \geq 
        \sup_{\substack{a \in \R \\ \xi \in H^1_0(\Omega) \\ a\geq 0}}
        \inf_{\substack{ %% begin constraints
        \psi \in \cal{H}
        }}\; %% end constraints
        \avg{
            a |\nabla \psi|^2 
            - \frac14 \abs{\nabla \xi}^2 
            - \frac{\Pe^2}{4} \abs{\nabla \invlap(\nabla \xi \cdot \nabla^\perp \psi)}^2 
            + f\xi
        } 
        - a.
    \end{align*}

    We now remove the term involving the inverse Laplacian using the identity 
    \begin{equation*}
        -\avg{ \frac{\Pe^2}{4} \abs{\nabla \invlap(\nabla \xi \cdot \nabla^\perp \psi)}^2 } = \min_{\theta \in H^1_0(\Omega)} \avg{ \abs{\nabla\theta}^2 - \Pe\,\xi\,\nabla^\perp \psi \cdot \nabla\theta },
    \end{equation*}
    arriving at the lower bound
    \begin{align*}
        \ocpval(\Pe)
        \geq 
        \sup_{\substack{a \in \R \\ \xi \in H^1_0(\Omega) \\ a\geq 0}}
        \inf_{\substack{ %% begin constraints
        \psi \in \cal{H} \\ 
        \theta \in H^1_0(\Omega)
        }}\; %% end constraints
        \avg{a\abs{\nabla\psi}^2 + \abs{\nabla\theta}^2 - \Pe\,\xi\,\nabla^\perp\psi \cdot \nabla\theta - \frac14 \abs{\nabla \xi}^2 + f\xi} - a.
    \end{align*}
    The minimax problem on the right-hand side of this inequality is equivalent to the maximization problem \cref{e:ldp} because the inner minimum is zero if the quadratic form $\QF(a,\xi)$ defined in \cref{e:qf} is positive semidefinite, and is negative infinity otherwise. Therefore, the outer maximization selects pairs $(a,\xi)$ satisfying $\QF(a,\xi)\succeq 0$. Since this constraint automatically implies that $a\geq 0$, the latter condition need not be enforced in \cref{e:ldp} and we conclude that $\ocpval(\Pe) \geq \lb(\Pe)$. Finally, $\lb(\Pe)\geq 0$ because the choices $a=0$ and $\xi=0$ are feasible for \cref{e:ldp} and have a zero objective value.
\end{proof}

\end{document}
